\section{SCENARIO \#13: "Fisher's Hill': September 22-26, 1864
}


\begin{scenariodata}
\textbf{Game Length:} 5 Turns\\
\end{scenariodata}

\begin{scenariodata}
(Weather: July - October Column)\\
\end{scenariodata}

\begin{scenariodata}
\textbf{First Phase:} Union\\
\end{scenariodata}

\begin{scenariodata}
\textbf{Union:}\\
\end{scenariodata}

\begin{scenariodata}
Resource Factor 215\\
\end{scenariodata}

\begin{scenariodata}
\textbf{Initial Placement:}\\
\end{scenariodata}

\begin{scenariodata}
Hex A-16 (Romney) > (41, 1A], (1F, 15)\\
Hex N-24 : 651, 11A, 9C, 3H,\\
\textbf{and adjacent :} (Sheridan, Wright,\\
\end{scenariodata}

Emory, Crook, 35, 3W)

\begin{scenariodata}
Hex \$-28 (Front Royal) : 9C, 3H,\\
\end{scenariodata}

\begin{scenariodata}
\textbf{and adjacent :} (1S, 1W)\\
Hex N-20 (Winchester) : 51, (1F, 28, 5W)\\
B\&O : (201, SA, 2C], (4F, 45)\\
\end{scenariodata}

\begin{scenariodata}
\textbf{Withdrawal Hexes:}\\
\end{scenariodata}

\begin{scenariodata}
N-7, 5-10, 2-13, CC-15 through CC-33\\
\end{scenariodata}

\begin{scenariodata}
(Supply/Wagon Entry Hexes:\\
\end{scenariodata}

\begin{scenariodata}
N-7, \$-10, Z-13, A-11, A-13, A-16, CC-15 through CC-33)\\
\end{scenariodata}

\begin{scenariodata}
One additional Fortification counter available.\\
\end{scenariodata}

\begin{scenariodata}
\textbf{Confederates:}\\
\end{scenariodata}

\begin{scenariodata}
Resource Factor 211\\
\end{scenariodata}

\begin{scenariodata}
\textbf{Initial Placement:}\\
\end{scenariodata}

\begin{scenariodata}
Hex R-31 7\\
\end{scenariodata}

\begin{scenariodata}
\textbf{and adjacent :} SC, 2H, (1S, 1W)\\
\end{scenariodata}

\begin{scenariodata}
Hexes L-26, M-27, N-28 :\\
\end{scenariodata}

\begin{scenariodata}
\textbf{and adjacent :} 351, 8A, 5C, 2H,\\
(Early, Gordon,\\
\end{scenariodata}

Ramseur, 4F, 2S, 4W)

\begin{scenariodata}
Hex T-33 and adjacent :1\\
\end{scenariodata}

\begin{scenariodata}
Hex 0-51 (Staunton) : UI, (1F, 1S)\\
\end{scenariodata}

\begin{scenariodata}
Hex BB-53 (Charlottesville) : 21, 1A, (1F)\\
\end{scenariodata}

\begin{scenariodata}
(Partisans: McNeill)\\
\end{scenariodata}

\begin{scenariodata}
Hexes CC-40 through CC-59 : Available Turn \#5\\
\end{scenariodata}

\begin{scenariodata}
Adds + 2 Resource Factors : 101, 4A, (Kershaw)\\
\end{scenariodata}

\begin{scenariodata}
\textbf{Withdrawal Hexes:}\\
\end{scenariodata}

\begin{scenariodata}
CC-40 through CC-58\\
\end{scenariodata}

\begin{scenariodata}
(Supply/Wagon Entry Hexes:\\
\end{scenariodata}

\begin{scenariodata}
CC-40 through CC-58, 0-51, M-50, E-46)\\
\end{scenariodata}

\begin{scenariodata}
No additional Fortification counters.\\
\end{scenariodata}

\begin{scenariodata}
(SPECIAL: Devastated areas are shown by placing Destruction counters\\
\end{scenariodata}

in the following hexes: R-23, P-24, O-21, M-22, L-19.)

\begin{scenariodata}
\textbf{Victory Conditions:} All other results are a draw.\\
\end{scenariodata}

\begin{scenariodata}
\textbf{ADVANCED GAME:} Union wins if they have 20 or more lead in total\\
\end{scenariodata}

\subsection*{Victory Points.}

\begin{scenariodata}
Confederates win if they have 20 or more lead in\\
\end{scenariodata}

total Victory Points.

\begin{scenariodata}
\textbf{BASIC GAME:} Same as the Advanced Game, only "spot'' the Con-\\
\end{scenariodata}

federates 15 Victory Points at the start of the game.

\begin{scenariodata}
\textbf{HISTORICAL RESULTS:} Union Victory.\\
\end{scenariodata}

\begin{scenariodata}
\textbf{PLAYTESTER'S COMMENTS:} The Union has all the strength necessary\\
\end{scenariodata}

to force the Confederate line, and head south. The Confederates should avoid being cut off, try to fight a delaying action with.as few casualties as possible, and finally break contact completely.

A flank attack by Crook's VIII Corps totally unhinged the over- extended Confederate line, and the defeated Confederates were again sent "'a whirling" through the Valley. The Union forces followed, and Early proved too weak to stop them. The Virginia Central Railroad was badly damaged (although quickly repaired by the Confederates), and the fertile area around Staunton devastated. Early had obtained a few reinforcements, but continued to maneuver elusively, and avoid battle. Orders to finish up his business in the valley, and transfer his units elsewhere led Sheridan to decide that it was time to return north...

